\section{Discusión}

Los resultados sugieren dos conclusiones prácticas para el diseño de
asistentes domóticos locales. En primer lugar, el tamaño del modelo mejora
de forma clara la interpretación de qué dispositivo, qué acción y en qué
ubicación (intent, dispositivo, ubicación), pero el costo en latencia crece
más rápido que la ganancia en exactitud: pasar de Qwen2.5-0.5B a
SmolLM2-1.7B multiplica la latencia por algo más de 4 (14,7s → 61,9s) a
cambio de una mejora de 15,6 puntos porcentuales en coincidencia exacta
(43,8\% → 59,4\%). Para un hub domótico real, donde la experiencia de
usuario depende de una respuesta rápida, Qwen2.5-1.5B-Instruct ofrece el
mejor compromiso observado en este estudio: una coincidencia exacta similar
a la del modelo más grande (50,0\% vs. 59,4\%) con una latencia bastante
menor (43,2s vs. 61,9s), y la mejor exactitud de campo ``valor'' (90,6\%) de
los cuatro modelos.

En segundo lugar, y quizás más relevante, el estancamiento en los campos
valor/unidad y el desempeño uniformemente bajo en la categoría de ajustes
relativos indican que escalar el modelo dentro del rango evaluado (0,36B a
1,7B) no alcanza para resolver esta clase de comando. Esto tiene una
implicancia de diseño concreta: un sistema domótico real podría
beneficiarse de una capa de post-procesamiento basada en reglas (por
ejemplo, detectar la ausencia de un número en el comando original y forzar
valor=null y unidad=null en la salida del modelo) en lugar de depender de
que el modelo aprenda esta restricción únicamente a partir de mayor escala.

El resultado llamativo de que SmolLM2-360M-Instruct obtuvo la mayor
exactitud en la categoría de consulta de estado (66,7\% vs. 33,3\% en los
otros tres) es probablemente ruido estadístico dado el tamaño reducido de
esa categoría (n=3) y no debería interpretarse como que el modelo más chico
interpreta mejor las consultas; un dataset más amplio por categoría es
necesario para confirmar o descartar este patrón.
